\section{Tools}\label{sec:tools}

\subsection{Regular variation}\label{sec:rv}

We use the following facts about a function $T$ that is regularly varying with
index $-\alpha\in(-1,0)$, all from~\cite[Sections~1.2 to~1.5]{bgt1989}.
\begin{enumerate}[(R1)]
\item (Uniform convergence.) $T(\lambda x)/T(x)\to\lambda^{-\alpha}$ uniformly for
$\lambda$ in compact subsets of $(0,\infty)$.
\item (Potter.) For every $c\ge1$ there is $C_c$ with $T(x/c)\le C_cT(x)$ for $x\ge c$.
\item (Karamata.) $\sum_{n\le w}T(n)\sim wT(w)/(1-\alpha)$ and
$\sum_{1\le i\le w}1/(iT(i))\sim1/(\alpha T(w))$ as $w\to\infty$.
\item (Dyadic sums.) $\sum_{l\ge1,\,2^l\le x}2^lT(2^{l-1})\le CxT(x)$ and
$\sum_{l\ge1,\,2^l>x}T(2^{l-1})\le CT(x)$ for $x\ge1$.
\end{enumerate}
(R4) follows from (R2) and Potter's theorem~\cite[Theorem~1.5.6]{bgt1989}, since
$wT(w)$ is regularly varying with the positive index $1-\alpha$ and $T$ with the
negative index $-\alpha$, so both sums are dominated by geometric series
(Lemma~\ref{lem:rv-facts}). Under (BAL) the two tails are comparable,
$c\,T_\plus\le T_\minus\le C\,T_\plus$, and $mT_\plus(a_m)\le C$ for all $m\ge1$. We
also use $a_{\lceil m/3\rceil}\ge c\,a_m$, which follows from (R2) for the function
$a$.

\subsection{Cited theorems}\label{sec:cited}

For a sign with $\alpha_\sgn<1$, under (RV) and (SPAN), $S^\sgn_m/a_\sgn(m)$ converges in law to
the positive stable variable with L\'evy tail $z^{-\alpha_\sgn}$, whose density we write
$g_\sgn$. Under (BAL) we have $mT_\minus(a_mz)\to bz^{-\alpha}$, so with the common norming $a_m$
the minus sums converge to $b^{1/\alpha}Z$, whose density is
$g_b(z)=b^{-1/\alpha}g(zb^{-1/\alpha})$. Note that $t^{-1/\alpha}g_b(yt^{-1/\alpha})=p_{bt}(y)$.
The next lemma is stated for one sign and applies to each sign whose index is in
$(0,1)$.

\begin{lemma}[classical inputs]\label{lem:cd}
Let $T=T_\sgn$ be regularly varying with index $-\alpha\in(-1,0)$ and let $f=f_\sgn$ satisfy
(SPAN). Write $u_m$, $u$ and $a_m$ for the objects of this sign.
\begin{enumerate}[(i)]
\item (Gnedenko's local limit theorem~\cite[Theorem~4.2.1]{ibragimovlinnik1971}.)
$\sup_x|a_mu_m(x)-g_\sgn(x/a_m)|\to0$ as $m\to\infty$. Hence $\sup_xu_m(x)\le C/a_m$ for all
$m\ge1$. Under (BAL) the same holds for the minus sums with the norming $a_m$ and the
density $g_b$.
\item (Local large deviations~\cite[Theorem~1.1, Eq.~(1.6)]{caravennadoney2019}.)
There is $C$ with $u_m(x)\le C\,\dfrac m{a_m}\,T(x)$ for all $m\ge1$ and $x\ge1$.
\item (The strong renewal theorem.) If $\alpha>1/2$, (SRT$_\sgn$)
holds~\cite{garsialamperti1963,erickson1970}, see~\cite[Theorem~1.4]{caravennadoney2019}.
If $\alpha\le1/2$, (SRT$_\sgn$) holds if and only if
\[
 \lim_{\delta\to0}\limsup_{x\to\infty}\,xT(x)\sum_{1\le z\le\delta x}
 \frac{f(x-z)}{zT(z)^2}=0
\]
\cite[Theorem~1.4]{caravennadoney2019}. In every case (LB$_\sgn$) implies
(SRT$_\sgn$)~\cite{doney1997}, see~\cite[Eq.~(1.15)]{caravennadoney2019}.
\item (The strong renewal theorem as a small-$m$ condition~\cite[Section~4.1,
Eq.~(4.1)]{caravennadoney2019}.) (SRT$_\sgn$) holds if and only if
\begin{equation}\label{eq:cd41}
 \lim_{\delta\to0}\limsup_{x\to\infty}I_\sgn(\delta,x)=0,\qquad
 I_\sgn(\delta,x)=xT(x)\sum_{1\le m\le\delta/T(x)}u_m(x).
\end{equation}
\item (Doney's local bound~\cite[Theorem~2.4]{berger2019ptrf}, extending~\cite[Theorem~A]{doney1997}.)
Under (LB$_\sgn$) there is $C$ with $u_m(x)\le C\,mT(x)/x$ for all $m\ge1$
and $x\ge a_m$.
\end{enumerate}
\end{lemma}

We comment on the form of these statements. In (ii) and (iv) Caravenna and Doney
use the tail $\Prob(X>x)$ and a function $A$ with $A(x)\sim1/\Prob(X>x)$ for renewal
laws. The difference between $\Prob(X>x)$ and $T(x)=\Prob(X\ge x)$ is
$f(x)=o(T(x))$. In (iv) their condition reads
$\lim_{\delta\to0}\limsup_x(x/A(x))\sum_{1\le m\le A(\delta x)}u_m(x)=0$. Since
$A(\delta x)\sim\delta^\alpha/T(x)$ and the sum increases with its upper limit, this
is equivalent to~\eqref{eq:cd41}. (SPAN) makes the span $1$, so their
$\Prob(S_m\in x+(-1,0])$ is $u_m(x)$. In (v) Berger uses a scale $a'_m$ that may differ
from ours by a constant factor. For $c\,a_m\le x\le a_m$ the bound follows from (ii),
since $m/a_m\le m/x$ there, so (v) holds for $x\ge a_m$ with our scale. For a sign with
index in $(1,2)$ we use the local limit theorem and the local large deviation bound in
the form of~\cite[Lemma~A.1]{paperF}: with $\mu=\E X$ and the scale $a$ of that sign,
$\sup_xu_m(x)\le C/a(m)$, and $u_m(x)\le C(m/a(m))T(x-m\mu)$ for $x-m\mu\ge1$.

\subsection{Elementary bounds}\label{sec:elementary}

The following bounds use only (RV) with index in $(0,1)$ and (SPAN) for one sign, and
we drop the sign from the notation. They hold for each sign with $\alpha_\sgn<1$.

\begin{lemma}[sums of runs]\label{lem:sums}
\begin{enumerate}[(i)]
\item $\Prob(S_m\ge x)\le C\,mT(x)$ for all $m,x\ge1$.
\item $\Prob(S_m<x)\le e^{-mT(x)}$ for all $m\ge0$ and $x\ge1$.
\item $u_m(x)\le C\,a_m^{-1}e^{-mT(x+1)/3}$ for all $m\ge2$ and $x\ge0$.
\end{enumerate}
\end{lemma}

\begin{proof}
(i) If $S_m\ge x$, then some $X_i\ge x$, or $\sum_iX_i\ones_{\{X_i<x\}}\ge x$. By the
union bound and Markov's inequality,
$\Prob(S_m\ge x)\le mT(x)+m\,\E[X;X<x]/x$. Now
$\E[X;X<x]=\sum_{1\le n<x}\Prob(n\le X<x)\le\sum_{n<x}T(n)\le CxT(x)$ by (R3).

(ii) $S_m<x$ forces $X_i<x$ for every $i\le m$, which has probability
$(1-T(x))^m\le e^{-mT(x)}$.

(iii) Let $h=\lfloor m/2\rfloor\ge m/3$. Conditioning on $S_h$,
$u_m(x)=\sum_{y\le x}u_h(y)u_{m-h}(x-y)\le\Prob(S_h\le x)\sup_zu_{m-h}(z)$. By (ii)
$\Prob(S_h\le x)=\Prob(S_h<x+1)\le e^{-mT(x+1)/3}$, and by Lemma~\ref{lem:cd}(i)
$\sup_zu_{m-h}(z)\le C/a_{m-h}\le C/a_m$, since $m-h\ge m/2$.
\end{proof}

The crude bound (i) has a sharp form on the range where one big summand carries the
sum. We need it uniformly in the number of summands.

\begin{lemma}[one big jump]\label{lem:onejump}
For $\varepsilon>0$ and $x\ge1$ put
$\varrho(\varepsilon,x)=\sup\bigl\{|\Prob(S_m\ge y)/(mT(y))-1|:\ y\ge x/2,\ m\ge1,\ mT(y)\le\varepsilon\bigr\}$.
Then $\lim_{\varepsilon\to0}\limsup_{x\to\infty}\varrho(\varepsilon,x)=0$.
\end{lemma}

\begin{proof}
Let $y\ge x/2$ and $mT(y)\le\varepsilon$. For the lower bound,
$\Prob(S_m\ge y)\ge\Prob(\max_{i\le m}X_i\ge y)=1-(1-T(y))^m\ge mT(y)-\tfrac12(mT(y))^2\ge mT(y)(1-\varepsilon/2)$.
For the upper bound fix $0<\delta<\eta<\frac12$ with $\eta/\delta\ge2$. On
$\{S_m\ge y\}$ one of the following holds: some $X_i\ge(1-\eta)y$; at least two summands
are at least $\delta y$; exactly one summand is at least $\delta y$, it is below $(1-\eta)y$,
and the sum $W$ of the summands below $\delta y$ is at least $\eta y$; no summand is at least
$\delta y$ and $W\ge y$. The first event has probability at most
$mT((1-\eta)y)\le mT(y)((1-\eta)^{-\alpha}+o(1))$ by (R1), the $o(1)$ being uniform in $m$ and
$y\ge x/2$ as $x\to\infty$. The second has probability at most $m^2T(\delta y)^2\le C_\delta\varepsilon\,mT(y)$
by (R2). For the last two we bound $W$ by a Chernoff estimate. Each summand of $W$ is
at most $\delta y$, so $e^{\lambda X}\le1+(e^{\lambda\delta y}-1)X/(\delta y)$ on $\{X<\delta y\}$,
and with $\E[X;X<\delta y]\le C\delta yT(\delta y)\le C_\delta\,\delta\,yT(y)$ by (R3) and (R2),
\[
 \Prob(W\ge z)\le\exp\Bigl(-\lambda z+m\,\frac{e^{\lambda\delta y}-1}{\delta y}\,C_\delta\delta yT(y)\Bigr)
 \le\Bigl(\frac{eC_\delta\delta\,mT(y)\,y}{z}\Bigr)^{z/(\delta y)}
\]
for $z\ge eC_\delta\delta\,mT(y)y$, by the choice $e^{\lambda\delta y}=z/(C_\delta\delta mT(y)y)$. At
$z=\eta y$ the condition holds for small $\varepsilon$, the base is $C_{\delta,\eta}mT(y)\le1$ and the
exponent is $\eta/\delta\ge2$, so
$\Prob(W\ge\eta y)\le(C_{\delta,\eta}mT(y))^2\le C_{\delta,\eta}\varepsilon\,mT(y)$. The third event has probability at most
$mT(\delta y)\Prob(W'\ge\eta y)$ with $W'$ a sum of $m-1$ truncated summands, which is at most
$C_\delta mT(y)\cdot C_{\delta,\eta}\varepsilon$. The fourth has probability at most
$\Prob(W\ge y)\le\Prob(W\ge\eta y)$. Adding up,
$\Prob(S_m\ge y)\le mT(y)[(1-\eta)^{-\alpha}+o(1)+C_{\delta,\eta}\varepsilon]$. Let $x\to\infty$, then
$\varepsilon\to0$, then $\eta\to0$.
\end{proof}

\begin{lemma}[the covering kernel at scale $N$]\label{lem:kernel-limit}
Assume (BAL) and (SPAN), and fix $\delta\in(0,1)$. Uniformly for $t\in[\delta,1/\delta]$ and
$x\in[\delta,1]$,
\[
 q^\plus_{m_N(t)+1}(\lfloor xN\rfloor)=T_\plus(N)\bigl[Q_t(x)+o(1)\bigr],\qquad
 q^\minus_{m_N(t)}(\lfloor xN\rfloor)=T_\plus(N)\bigl[Q^b_t(x)+o(1)\bigr].
\]
\end{lemma}

\begin{proof}
Write $m=m_N(t)$ and $k=\lfloor xN\rfloor$, and fix $\theta\in(0,\delta/2)$. Split
$q^\plus_{m+1}(k)=\sum_{r<k}u^\plus_m(r)T_\plus(k-r)$ at $r=k-\theta N$. On
$r\le k-\theta N$ we have $(k-r)/N\in[\theta,1]$, so by (R1)
$T_\plus(k-r)/T_\plus(N)=((k-r)/N)^{-\alpha}+o(1)$ uniformly. By
Lemma~\ref{lem:cd}(i), $u^\plus_m(r)=a_m^{-1}[g(r/a_m)+o(1)]$ uniformly in $r$. The error
terms contribute at most $o(1/a_m)\cdot k\cdot T_\plus(\theta N)=o(T_\plus(N))$, since
$N/a_m$ is bounded. The main term is
$T_\plus(N)\sum_{r\le k-\theta N}a_m^{-1}g(r/a_m)((k-r)/N)^{-\alpha}$, a Riemann sum with mesh
$1/N$ in $z=r/N$. Since $a_m/N\to t^{1/\alpha}$ uniformly and $g$ is bounded and uniformly
continuous, it converges to $\int_0^{x-\theta}p_t(z)(x-z)^{-\alpha}dz$, uniformly in
$t$ and $x$. On $r>k-\theta N$, Lemma~\ref{lem:cd}(i) and (R3) give at most
$(C/a_m)\sum_{i<\theta N+1}T_\plus(i)\le C_\delta\,\theta^{(1-\alpha)/2}T_\plus(N)$ for large
$N$, using (R1) and (R2) to compare $\theta NT_\plus(\theta N)$ with $NT_\plus(N)$. By the
first bound in~\eqref{eq:p-bound} the integral over $[x-\theta,x]$ is at most
$C_\delta\theta^{1-\alpha}$. Let $N\to\infty$ and then $\theta\to0$. The minus kernel is the
same with $u^\minus_{m-1}$ in place of $u^\plus_m$, the density $g_b$ in place of $g$,
and $T_\minus(N)\sim bT_\plus(N)$. The shift from $m$ to $m-1$ does not change the
limit.
\end{proof}

The next bound is the key estimate of the paper. It bounds the covering kernel by
the tail at the level, uniformly in the number of runs, under regular variation
alone.

\begin{lemma}[the key bound]\label{lem:key}
For a sign with $\alpha_\sgn<1$ there is $C$ with $q^\sgn_{m+1}(k)\le C\,T_\sgn(k)$ for all
$m\ge0$ and $k\ge1$.
\end{lemma}

\begin{proof}
We drop the sign. For $m=0$, $q_1=T$. Let $m\ge1$. By~\eqref{eq:q-cover},
$q_{m+1}(k)=\sum_{r<k}u_m(r)T(k-r)$. The part with $r<k/2$ is at most
$T(k/2)\le CT(k)$ by (R2). For $r\ge k/2$ put $i=k-r\in[1,k/2]$, and split the range
of $i$ into dyadic blocks $[w/2,w)$ with $w=2^l$, $l\ge1$, $w\le k$. On a block,
\[
 \sum_{w/2\le i<w}u_m(k-i)T(i)\le T(w/2)\,\Prob\bigl(S_m\in(k-w,k-w/2]\bigr).
\]
The window has at most $w$ lattice points, all at least $k/2$. By
Lemma~\ref{lem:cd}(ii) and (R2) each point has probability at most
$C(m/a_m)T(k)$. By Lemma~\ref{lem:sums}(i) and (R2) the whole window has probability
at most $C\,mT(k)$. So the block is at most
$C\,mT(k)\,T(w/2)\min\{1,w/a_m\}$. Summing with (R4), the blocks with $w\le a_m$ give
at most $C\frac{mT(k)}{a_m}\sum_{w\le a_m}wT(w/2)\le C\,mT(k)T(a_m)$, and the blocks
with $w>a_m$ give at most $C\,mT(k)\sum_{w>a_m}T(w/2)\le C\,mT(k)T(a_m)$. Since
$mT(a_m)\le C$, the total is at most $CT(k)$.
\end{proof}

Note that Lemma~\ref{lem:cd}(ii) is used here at points $x\ge k/2$ that may lie far
below $a_m$ as well as above it, and it holds for all $x$.

\begin{lemma}[decay in $m$]\label{lem:decay}
For a sign with $\alpha_\sgn<1$ there are $C,c$ with $q^\sgn_m(j)\le C\,T_\sgn(j)\,e^{-c\,mT_\sgn(j)}$ for
all $m,j\ge1$.
\end{lemma}

\begin{proof}
We drop the sign. If $mT(j)\le9$, this is Lemma~\ref{lem:key} with $c=1/9$. Let
$mT(j)>9$, so $m\ge10$, and put $h=\lfloor(m-1)/2\rfloor\ge m/3$. Splitting
$S_{m-1}=S_h+(S_{m-1}-S_h)$ gives $q_m(j)=\sum_{s<j}u_h(s)\,q_{m-h}(j-s)$. For $s<j/2$
we use $q_{m-h}(j-s)\le CT(j-s)\le CT(j)$ (Lemma~\ref{lem:key} and (R2)) and
$\Prob(S_h<j/2)\le e^{-hT(j)}\le e^{-mT(j)/3}$ (Lemma~\ref{lem:sums}(ii)). For
$j/2\le s<j$ we use
\[
 \sum_{j/2\le s<j}u_h(s)\,q_{m-h}(j-s)\le\sup_{s<j}u_h(s)\sum_{i\le j/2}CT(i)
 \le C\,a_h^{-1}e^{-hT(j)/3}\,jT(j),
\]
by Lemma~\ref{lem:sums}(iii) and (R3). Since $hT(j)\ge mT(j)/3>3$, we have
$T(j)\ge3/h\ge T(a_h)$ for large $h$, so $j\le Ca_h$ and $jT(j)/a_h\le CT(j)$. Both
parts are at most $CT(j)e^{-mT(j)/9}$.
\end{proof}

The next lemma collects the Fourier bounds we need for the differences of the convolution
powers, in space and in the index.

\begin{lemma}[Fourier bounds]\label{lem:fourier}
Let $\alpha_\sgn<1$, assume (RV) and (SPAN) for the sign $\sgn$, and drop the sign. Let
$\varphi(t)=\E e^{itX}$, and for $y\ge1$ let $V^{(y)}_m(x)=\Prob(S_m=x,\ \max_{i\le m}X_i<y)$.
\begin{enumerate}[(i)]
\item There is $c>0$ with $|\varphi(t)|\le e^{-cT(1/|t|)}$ for $0<|t|\le\pi$.
\item For all $m\ge1$, $y\ge1$ and $x$, $|V^{(y)}_m(x-1)-V^{(y)}_m(x)|\le Ce^{CmT(y)}a_m^{-2}$. In
particular $\sup_x|u_m(x-1)-u_m(x)|\le Ca_m^{-2}$, and $|u_m(x-1)-u_m(x)|\le Ca_m^{-2}e^{-cmT(x)}$ for all
$m\ge1$ and $x\ge1$.
\item $\sup_x|u_{m+1}(x)-u_m(x)|\le C/(ma_m)$ for all $m\ge1$.
\end{enumerate}
\end{lemma}

\begin{proof}
(i) Let $X'$ be an independent copy of $X$. Then $|\varphi(t)|^2=\E\cos(t(X-X'))$, and for
$|t|\le1$, using $1-\cos\vartheta\ge\vartheta^2/4$ on $|\vartheta|\le1$, a point $n_0$ of the support of
$f$, and Karamata's theorem,
\[
 1-|\varphi(t)|^2\ge\frac{t^2}4\E\bigl[(X-X')^2;\ |X-X'|\le1/|t|\bigr]
 \ge\frac{t^2}4f(n_0)\sum_{2n_0\le n\le1/(2|t|)}\frac{n^2}4f(n)\ge c\,T(1/|t|)
\]
for $|t|\le t_0$. Hence $|\varphi(t)|\le(1-cT(1/|t|))^{1/2}\le e^{-cT(1/|t|)/2}$ there. On
$t_0\le|t|\le\pi$, (SPAN) gives $|\varphi(t)|<1$, so by continuity $|\varphi(t)|\le1-c_0$, which is at
most $e^{-cT(1/|t|)}$ for $c$ small since $T\le1$.

(ii) The generating function of $V^{(y)}_m$ is $\varphi_y^m$ with $\varphi_y(t)=\sum_{n<y}f(n)e^{itn}$, and
$|\varphi_y|\le|\varphi|+T(y)\le e^{-cT(1/|t|)}(1+e^cT(y))$, so $|\varphi_y(t)|^m\le e^{CmT(y)}e^{-cmT(1/|t|)}$.
Fourier inversion gives
\[
 |V^{(y)}_m(x-1)-V^{(y)}_m(x)|=\Bigl|\frac1{2\pi}\int_{-\pi}^{\pi}\varphi_y(t)^m(e^{it}-1)e^{-itx}\,dt\Bigr|
 \le\frac{e^{CmT(y)}}{2\pi}\int_{-\pi}^{\pi}|t|\,e^{-cmT(1/|t|)}\,dt .
\]
Substitute $t=s/a_m$ and fix $\delta\in(0,\alpha)$. For $s\le1$ the exponential is at most $1$.
For $1\le s\le a_m/x_0$, Potter's bounds give $mT(a_m/s)\ge cs^{\alpha-\delta}$. For $|t|\ge1/x_0$ the integrand is at most
$e^{-cmT(x_0)}\le e^{-c'm}$. So the integral is at most
$a_m^{-2}\bigl[\int_0^1s\,ds+\int_1^\infty se^{-cs^{\alpha-\delta}}ds\bigr]+Ce^{-c'm}\le Ca_m^{-2}$, since
$a_m\le Cm^{1/\alpha+\delta}$. The case $y=\infty$ is the bound for $u_m$. For the last claim let
$m\ge2$, $k=\lceil m/2\rceil$ and write $u_m=u_{m-k}*u_k$. Then
$|u_m(x-1)-u_m(x)|\le\sup_z|u_k(z-1)-u_k(z)|\,\Prob(S_{m-k}\le x)\le Ca_k^{-2}e^{-(m-k)T(x+1)}$ by
Lemma~\ref{lem:sums}(ii), and $a_k\ge ca_m$, $(m-k)T(x+1)\ge cmT(x)$.

(iii) We have $u_{m+1}(x)-u_m(x)=(2\pi)^{-1}\int_{-\pi}^\pi\varphi(t)^m(\varphi(t)-1)e^{-itx}dt$ and
$|\varphi(t)-1|\le\E\min(2,|t|X)\le2T(1/|t|)+|t|\,\E[X;X<1/|t|]\le CT(1/|t|)$ by Karamata's theorem. So
the left side is at most $C\int_0^\pi T(1/t)e^{-cmT(1/t)}dt$. Substitute $t=s/a_m$ again. For $s\le1$,
$T(a_m/s)\le T(a_m)\le C/m$. For $1\le s\le a_m/x_0$, $T(a_m/s)\le Cs^{\alpha+\delta}/m$ and the
exponential is at most $e^{-cs^{\alpha-\delta}}$. For $s\ge a_m/x_0$ the integrand is at most
$e^{-c'm}$ on an interval of length $\pi a_m$. Hence the integral is at most
$a_m^{-1}[C/m+C/m+Ca_me^{-c'm}]\le C/(ma_m)$.
\end{proof}

The last lemma of this section is the local form of the inverse renewal count. For
a sign with $\alpha_\sgn<1$ let $E_\sgn=Z_\sgn^{-\alpha_\sgn}$, where $Z_\sgn$ is the positive
stable variable with L\'evy tail $z^{-\alpha_\sgn}$, and let $\phi_\sgn$ be its density. Since
$\Prob(T_\sgn(n)\#\{m:S^\sgn_m<n\}\ge t)=\Prob(S^\sgn_{\lfloor t/T_\sgn(n)\rfloor}<n)\to\Prob(Z_\sgn<t^{-1/\alpha_\sgn})$,
the variable $E_\sgn$ is the limit of $T_\sgn(n)$ times the number of runs of sign $\sgn$
completed within $n$ steps of that sign. It is a rescaled Mittag-Leffler variable, with
$\E E_\sgn^n=n!/(\Gamma(1+n\alpha_\sgn)\Gamma(1-\alpha_\sgn)^n)$ and $\phi_\sgn(0)=1$, because
$\Prob(E_\sgn\le\varepsilon)=\Prob(Z_\sgn\ge\varepsilon^{-1/\alpha_\sgn})\sim\varepsilon$. Its density is
decreasing from $0$ when $\alpha_\sgn<1/2$ and has an interior maximum when $\alpha_\sgn>1/2$,
since $\phi'_\sgn(0)=-\Gamma(1-\alpha_\sgn)^2/\Gamma(1-2\alpha_\sgn)$, and at $\alpha_\sgn=1/2$ it is
$\phi_\sgn(x)=e^{-\pi x^2/4}$.

\begin{lemma}[the local inverse renewal count]\label{lem:Q}
Let $\alpha_\sgn<1$, assume (RV) and (SPAN) for the sign $\sgn$, and drop the sign.
\begin{enumerate}[(i)]
\item $q_{m+1}(k)\ge(1-CmT(k))\,T(k)$ for all $m\ge0$ and $k\ge1$.
\item For every $K>0$, $\sup_{m\ge0:\ mT(k)\le K}\bigl|q_{m+1}(k)/T(k)-\phi(mT(k))\bigr|\to0$ as $k\to\infty$.
\end{enumerate}
\end{lemma}

\begin{proof}
(i) Since $T(k-r)\ge T(k)$ for $r\ge0$, \eqref{eq:q-cover} gives
$q_{m+1}(k)\ge\Prob(S_m<k)T(k)\ge(1-CmT(k))T(k)$ by Lemma~\ref{lem:sums}(i).

(ii) First let $t=mT(k)\in[\varepsilon,K]$. Then $m\to\infty$ with $k$, and
$a_m/k\to t^{1/\alpha}$ uniformly, as in Lemma~\ref{lem:rv-facts}(iii). Fix $\eta\in(0,1/2)$ and
split $q_{m+1}(k)=\sum_{r<k}u_m(r)T(k-r)$ at $r=(1-\eta)k$. As in the proof of
Lemma~\ref{lem:kernel-limit}, the part $r<(1-\eta)k$ is a Riemann sum that converges,
uniformly in $t\in[\varepsilon,K]$, to
$T(k)\int_0^{(1-\eta)t^{-1/\alpha}}g(w)(1-t^{1/\alpha}w)^{-\alpha}dw$, and the part
$r\ge(1-\eta)k$ is at most $(C/a_m)\sum_{i\le\eta k}T(i)\le C_\varepsilon\eta^{1-\alpha}T(k)$. Letting
$\eta\to0$ gives the limit $\int_0^1p_t(x)(1-x)^{-\alpha}dx$, and this is $\phi(t)$. Indeed, let
$(\sigma_s)_{s\ge0}$ be the subordinator with L\'evy measure $\alpha r^{-\alpha-1}dr$ and no drift, so
that $\sigma_t$ has the density $p_t$, and let $\tau=\inf\{s:\sigma_s>1\}$. The compensation
formula~\cite[Proposition~III.2]{bertoin1996} gives $\tau$ the density
$t\mapsto\int_0^1p_t(x)\,(1-x)^{-\alpha}\,dx$, since $(1-x)^{-\alpha}$ is the L\'evy tail at $1-x$.
Moreover $\{\tau>t\}=\{\sigma_t\le1\}=\{t^{1/\alpha}Z\le1\}=\{E\ge t\}$, so $\tau$ has the law of $E$.

Now let $t=mT(k)\le\varepsilon$. The lower bound (i) gives $q_{m+1}(k)/T(k)\ge1-C\varepsilon$, and
$\phi(t)\ge1-C\varepsilon$ on $[0,\varepsilon]$. For the upper bound fix $0<\theta,\eta<\frac12$ and
with $\eta<\frac14$, and write $\delta_k=\sup_{1/2\le\lambda\le1}|T(\lambda k)/T(k)-\lambda^{-\alpha}|$,
which tends to $0$ by (R1).
Split $q_{m+1}(k)$ into three ranges. The range $r<\eta k$ contributes at most
$T((1-\eta)k)\le((1-\eta)^{-\alpha}+\delta_k)T(k)$. The range $\eta k\le r\le(1-\eta)k$ contributes
at most $T(\eta k)\Prob(S_m\ge\eta k)\le C_\eta\varepsilon\,T(k)$ by Lemma~\ref{lem:sums}(i) and (R2).
For the range $r>(1-\eta)k$ put $i=k-r<\eta k$ and use dyadic blocks $[w/2,w)$ with $w<2\eta k$. A block
contributes at most $T(w/2)\Prob(S_m\in[k-w,k))$, and the window probability is at most
the smaller of $Cw(m/a_m)T(k)$, by Lemma~\ref{lem:cd}(ii), and
\begin{multline*}
 \Prob(S_m\ge k-w)-\Prob(S_m\ge k)
 \le mT(k)\Bigl[\frac{T(k-w)}{T(k)}-1+2\varrho(\varepsilon,k)\frac{T(k-w)}{T(k)}\Bigr]\\
 \le mT(k)\bigl[Cw/k+\delta_k+4\varrho(2\varepsilon,k)\bigr],
\end{multline*}
by Lemma~\ref{lem:onejump} at the levels $k-w\ge k/2$ and $k$, where $mT(k-w)\le2\varepsilon$, and
where we used $w<2\eta k<k/2$ and $T(k-w)/T(k)\le2$ for large $k$. The
blocks with $w\le\theta a_m$ contribute at most $C(mT(k)/a_m)\sum_{w\le\theta a_m}wT(w/2)\le C\theta^{1-\alpha'}mT(a_m)T(k)\le C\theta^{1-\alpha'}T(k)$
for any $\alpha'\in(\alpha,1)$, by (R4) and Potter's bounds. The blocks with
$\theta a_m<w<2\eta k$ contribute at most
$mT(k)\sum_{w}T(w/2)[Cw/k+\delta_k+4\varrho]\le CmT(k)[\eta T(\eta k)+(\delta_k+4\varrho)T(\theta a_m)]$ by (R4),
and $mT(\theta a_m)\le C\theta^{-\alpha'}$, so this is at most
$C[\eta^{1-\alpha}\varepsilon+(\delta_k+4\varrho(2\varepsilon,k))\theta^{-\alpha'}]T(k)$. Hence
\begin{multline*}
 \sup_{mT(k)\le\varepsilon}\frac{q_{m+1}(k)}{T(k)}
 \le(1-\eta)^{-\alpha}+\delta_k+C\theta^{1-\alpha'}+C_\eta\varepsilon\\
 +C\theta^{-\alpha'}\bigl(\delta_k+4\varrho(2\varepsilon,k)\bigr).
\end{multline*}
Let $k\to\infty$, then $\varepsilon\to0$, then $\theta\to0$, then $\eta\to0$. Together with the lower
bound and $\phi(t)=1+O(\varepsilon)$ on $[0,\varepsilon]$ this proves (ii) on $[0,\varepsilon]$, and with the
first part on $[\varepsilon,K]$.
\end{proof}
